
\documentclass[11pt]{article}

\usepackage[a4paper,margin=1in]{geometry}
\usepackage[T1]{fontenc}
\usepackage{lmodern}
\usepackage{amsmath,amssymb,amsthm}
\usepackage{booktabs}
\usepackage{longtable}
\usepackage{microtype}
\usepackage{hyperref}
\usepackage{xurl}

\hypersetup{
  colorlinks=true,
  linkcolor=black,
  citecolor=black,
  urlcolor=black,
  pdftitle={Outer-In Decimal Digit Permutations with Equal Prime Support},
  pdfauthor={[Author]}
}

\newtheorem{theorem}{Theorem}
\newtheorem{proposition}{Proposition}
\newtheorem{lemma}{Lemma}
\newcommand{\rad}{\operatorname{rad}}

\title{Outer-In Decimal Digit Permutations with Equal Prime Support}
\author{[Author]}
\date{22 September 2026}

\begin{document}
\maketitle

\begin{abstract}
We study the decimal digit transformation obtained by reading a number from the two ends inward. For
\[
n=d_1d_2\cdots d_k,
\]
define
\[
T(n)=d_1d_kd_2d_{k-1}d_3d_{k-2}\cdots .
\]
We seek positive integers \(n\) for which \(n\ne T(n)\) and \(n\) and \(T(n)\) have the same set of distinct prime divisors. The sequence considered here is restricted to integers not ending in \(0\). We prove an exact criterion in terms of \(g=\gcd(n,T(n))\) and the coprime reduced pair \(n/g,T(n)/g\). This gives a factorization-light exact test and leads to an exhaustive computation through \(10^{10}\). Exactly 22 integers satisfy the condition in the stated search domain. We also prove three infinite repeated-block families generated by \(224\mapsto242\), \(448\mapsto484\), and \(117\mapsto171\).
\end{abstract}

\section{Definition and motivation}

Let
\[
n=d_1d_2\cdots d_k
\]
be the usual decimal expansion of a positive integer, with \(d_1\ne0\). Define the outer-in transformation \(T\) by taking the first digit, then the last digit, then the second digit, then the second-to-last digit, and so on:
\[
T(n)=d_1d_kd_2d_{k-1}d_3d_{k-2}\cdots .
\]
For example,
\[
T(12345678)=18273645,\qquad
T(12345)=15243,\qquad
T(3024)=3402.
\]

For a positive integer \(n\), write
\[
\operatorname{supp}(n)=\{p:\ p\text{ is prime and }p\mid n\}
\]
and
\[
\rad(n)=\prod_{p\mid n}p.
\]
Thus \(\rad(n)=\rad(m)\) is equivalent to equality of the sets of distinct prime divisors of \(n\) and \(m\).

The sequence studied here consists of the integers satisfying
\[
n\ne T(n),\qquad \rad(n)=\rad(T(n)),
\tag{1}
\]
together with the explicit normalization
\[
n\not\equiv0\pmod {10}.
\tag{2}
\]
Condition (2) is part of the definition of the sequence; it is not needed in order for \(T(n)\) itself to be defined.

The condition \(n\ne T(n)\) excludes the fixed points of this particular digit permutation. It is not the same as requiring a number to be non-palindromic.

\section{An exact criterion}

\begin{theorem}
Let
\[
g=\gcd(n,T(n)),\qquad n=ag,\qquad T(n)=bg,
\]
where \(\gcd(a,b)=1\). Then
\[
\rad(n)=\rad(T(n))
\]
if and only if
\[
\rad(ab)\mid g.
\tag{3}
\]
Equivalently,
\[
a\,\rad(ab)\mid n.
\tag{4}
\]
\end{theorem}

\begin{proof}
Assume first that \(\rad(n)=\rad(T(n))\). If a prime \(p\) divides \(a\), then \(p\mid n\) and hence \(p\mid T(n)=bg\). Since \(\gcd(a,b)=1\), \(p\nmid b\), so \(p\mid g\). Interchanging \(a\) and \(b\) shows that every prime divisor of \(b\) also divides \(g\). Hence every prime dividing \(ab\) divides \(g\), proving (3).

Conversely, suppose \(\rad(ab)\mid g\). Every prime divisor of \(a\) or \(b\) then already divides \(g\). Consequently the prime divisors of \(ag=n\) are exactly the prime divisors of \(g\), and the same is true of \(bg=T(n)\). Therefore
\[
\rad(n)=\rad(T(n)).
\]
The equivalent form (4) follows from \(n=ag\).
\end{proof}

This criterion is the basis of the exhaustive search. It avoids independently factoring the two large integers \(n\) and \(T(n)\): after computing one gcd, it is enough to test whether every prime divisor of the two reduced factors already occurs in \(g\).

\section{Exhaustive computation}

The search constructs \(T(n)\) directly from the decimal digits. Before the exact test, necessary support checks are performed for
\[
2,5,7,11,13,17,19,23.
\]
A candidate is rejected if one of these primes divides exactly one of \(n\) and \(T(n)\). The primes \(3\) and \(9\) require no analogous test because a digit permutation preserves the digit sum.

Every surviving candidate is tested using Theorem 1:
\[
g=\gcd(n,T(n)),\qquad
a=n/g,\qquad
b=T(n)/g,
\]
followed by the exact divisibility condition \(\rad(ab)\mid g\).

The reported computation covers
\[
0<n<10^{10},\qquad n\not\equiv0\pmod {10}.
\]
The historical computation had two phases: an exhaustive search below \(10^9\), followed by a complete search of the ten-digit interval
\[
10^9\le n<10^{10}.
\]
The ten-digit phase used an audited 64-bit implementation. Since
\[
10^{10}<2^{64},
\]
all values in that interval are represented exactly by the implementation.

The resulting 22-term initial segment is
\[
\begin{aligned}
&3024,\ 6048,\ 10143,\ 20286,\ 27864,\ 117216,\ 189216,\ 224224,\\
&333234,\ 352872,\ 387585,\ 448448,\ 666468,\ 3892032,\ 3928824,\\
&137539323,\ 247153125,\ 321881616,\ 869148324,\ 1938285196,\\
&2240563518,\ 2952612351.
\end{aligned}
\]

\begin{proposition}
Among positive integers \(n<10^{10}\) with \(n\not\equiv0\pmod {10}\), exactly 22 integers satisfy (1).
\end{proposition}

The first 19 terms occur below \(10^9\); the complete ten-digit interval contains exactly three terms.

\begin{longtable}{r r r r}
\caption{The 22 verified terms, their images, gcds, and reduced ratios.}
\label{tab:data}\\
\toprule
$n$ & $T(n)$ & $\gcd(n,T(n))$ & $n/g:T(n)/g$\\
\midrule
\endfirsthead
\toprule
$n$ & $T(n)$ & $\gcd(n,T(n))$ & $n/g:T(n)/g$\\
\midrule
\endhead

3024 & 3402 & 378 & 4:9 \\
6048 & 6804 & 756 & 8:9 \\
10143 & 13041 & 1449 & 7:9 \\
20286 & 26082 & 2898 & 7:9 \\
27864 & 24768 & 3096 & 9:8 \\
117216 & 161172 & 14652 & 8:11 \\
189216 & 168192 & 21024 & 9:8 \\
224224 & 242242 & 2002 & 112:121 \\
333234 & 343332 & 10098 & 33:34 \\
352872 & 325728 & 27144 & 13:12 \\
387585 & 358875 & 14355 & 27:25 \\
448448 & 484484 & 4004 & 112:121 \\
666468 & 686664 & 20196 & 33:34 \\
3892032 & 3283902 & 121626 & 32:27 \\
3928824 & 3492288 & 436536 & 9:8 \\
137539323 & 133273593 & 60939 & 2257:2187 \\
247153125 & 254271135 & 395445 & 625:643 \\
321881616 & 362116818 & 40235202 & 8:9 \\
869148324 & 846293184 & 653004 & 1331:1296 \\
1938285196 & 1699318528 & 26551852 & 73:64 \\
2240563518 & 2821450356 & 82983834 & 27:34 \\
2952612351 & 2195532261 & 75708009 & 39:29 \\

\bottomrule
\endlastfoot
\end{longtable}

For the three ten-digit pairs, independent factorizations are
\[
1938285196=2^2\cdot73^2\cdot90931,
\qquad
1699318528=2^8\cdot73\cdot90931,
\]
\[
2240563518=2\cdot3^5\cdot17\cdot139\cdot1951,
\]
\[
2821450356=2^2\cdot3^2\cdot17^2\cdot139\cdot1951,
\]
and
\[
2952612351=3^3\cdot13^2\cdot29\cdot53\cdot421,
\]
\[
2195532261=3^2\cdot13\cdot29^2\cdot53\cdot421.
\]
Thus the equality of prime support in all three ten-digit cases can also be checked directly.

As a small example,
\[
3024=2^4\cdot3^3\cdot7,\qquad
3402=2\cdot3^5\cdot7,
\]
so both numbers have radical \(42\).

The computational completeness statement is only for the stated range and the nonzero-last-digit restriction.

\section{Repeated blocks and infinite families}

Let \(B\) and \(C\) be decimal blocks of the same length \(m\).

\begin{lemma}
If
\[
T(BB)=CC,
\tag{5}
\]
then for every \(r\ge1\),
\[
T(B^{2r})=C^{2r},
\tag{6}
\]
where \(B^{2r}\) denotes the decimal concatenation of \(2r\) copies of \(B\).
\end{lemma}

\begin{proof}
Write
\[
B=b_1b_2\cdots b_m.
\]
In the word \(B^{2r}\), consider the output in consecutive pairs. The first pair is \(b_1,b_m\), the second is \(b_2,b_{m-1}\), and so on through the \(m\)-digit block pattern
\[
b_1,b_m,b_2,b_{m-1},\ldots .
\]
Because the block \(B\) is repeated, exactly the same \(2m\)-digit pattern is produced for each of the \(r\) block-pairs. That pattern is \(T(BB)=CC\). Hence the complete output is \(r\) copies of \(CC\), namely \(C^{2r}\).
\end{proof}

Let \(A=[B]\) and \(C_0=[C]\) be the numerical values of the blocks, and define
\[
Q_r=1+10^m+10^{2m}+\cdots+10^{(2r-1)m}.
\]
Then
\[
[B^{2r}]=AQ_r,\qquad [C^{2r}]=C_0Q_r.
\tag{7}
\]
Moreover,
\[
10^m+1\mid Q_r,
\tag{8}
\]
because \(10^m\equiv-1\pmod{10^m+1}\) and \(Q_r\) contains an even number of powers.

Let
\[
g=\gcd(A,C_0),\qquad A=ag,\qquad C_0=bg,\qquad \gcd(a,b)=1.
\]
Combining (7), (8), and Theorem 1 gives the following sufficient condition.

\begin{proposition}
If
\[
\rad(ab)\mid g(10^m+1),
\tag{9}
\]
then every even repetition \(B^{2r}\) has the same prime support as the corresponding repetition \(C^{2r}\).
\end{proposition}

\begin{proof}
Since \(10^m+1\mid Q_r\), condition (9) implies
\[
\rad(ab)\mid gQ_r.
\]
Applying Theorem 1 to \(AQ_r\) and \(C_0Q_r\) proves equality of prime support for every \(r\ge1\).
\end{proof}

\subsection{\texorpdfstring{The family $224\mapsto242$}{The family 224 -> 242}}

We have
\[
224=2^5\cdot7,\qquad242=2\cdot11^2,
\]
so
\[
g=2,\qquad a=112,\qquad b=121,\qquad
\rad(ab)=2\cdot7\cdot11.
\]
For \(m=3\),
\[
10^3+1=1001=7\cdot11\cdot13.
\]
Hence (9) holds. Therefore, for every \(r\ge1\), the decimal integer obtained by repeating \(224\) exactly \(2r\) times is a solution, and its image under \(T\) is the corresponding repetition of \(242\).

The first member is
\[
224224\mapsto242242.
\]

\subsection{\texorpdfstring{The family $448\mapsto484$}{The family 448 -> 484}}

Similarly,
\[
448=2^6\cdot7,\qquad484=2^2\cdot11^2,
\]
so
\[
g=4,\qquad a=112,\qquad b=121.
\]
Again \(m=3\), and \(10^3+1\) contains both required support primes \(7\) and \(11\). Thus every even repetition of \(448\) is a solution and is mapped to the corresponding even repetition of \(484\).

\subsection{\texorpdfstring{The family $117\mapsto171$}{The family 117 -> 171}}

Take the 9-digit blocks
\[
B=117117117,\qquad C=171171171.
\]
Directly,
\[
T(BB)=CC,
\]
or
\[
T(117117117117117117)
=
171171171171171171.
\]
Also,
\[
\gcd(B,C)=9009009,\qquad
B/\gcd(B,C)=13,\qquad
C/\gcd(B,C)=19.
\]
Thus the reduced pair is \(13:19\). Finally,
\[
10^9+1=7\cdot11\cdot13\cdot19\cdot52579.
\]
Hence (9) holds. Every even repetition of \(B\), equivalently every repetition of \(117\) exactly \(6r\) times, is therefore a solution and is mapped to the corresponding repetition of \(171\).

The first member of this family is
\[
117117117117117117
\mapsto
171171171171171171.
\]

\section{Relation to previous work}

The closest OEIS analogue is A110751, which records integers whose decimal digit reversal has the same prime divisors; its non-palindromic subsequence is A110819 \cite{OEISA110751,OEISA110819}. The present transformation is different: it alternates digits from opposite ends rather than reversing the whole digit string.

A244514 gives another nearby example: it studies integers for which all cyclic permutations of the decimal digits have the same prime divisors \cite{OEISA244514}. The present problem fixes one specific digit permutation rather than requiring the property for a whole cyclic orbit. A371307 records the analogous smallest nonpalindromic examples for digit reversal in varying bases.

More generally, digit-permutation arithmetic is studied in the theory of permutiples. Holt's papers treat integers that are nontrivial multiples of permutations of their digits \cite{Holt2017,Holt2024}. No multiplicative relation such as \(T(n)=qn\) is imposed here; the invariant is equality of the sets of distinct prime divisors.

Searches of the OEIS and web-indexed sources performed for this work through 22 September 2026 did not locate an exact match for the displayed initial terms together with the outer-in/equal-prime-support definition. This is a report of the searches performed, not a formal priority claim.

\section{Reproducibility and limitations}

The accompanying archive separates the exploratory research history from a clean reproducibility release. The release contains the exhaustive-search source, an independent verifier, the consolidated dataset, the raw ten-digit outputs, the OEIS data, and verification records.

The search program operates on one decimal digit length at a time, so the width used to construct \(T(n)\) is fixed during each run. The ten-digit interval is therefore searched as
\[
[10^9,10^{10}).
\]
The implementation uses 64-bit unsigned integers, which are sufficient throughout this bounded domain.

The independent verifier reconstructs \(T(n)\), computes the gcd and reduced pair, checks the exact criterion of Theorem 1, and checks the expected image for each candidate. It also performs explicit factorization checks for the 22 listed pairs. Those factorizations are used as independent certificates for the present dataset, not as a general-purpose large-integer factoring method.

The exhaustive result should therefore be read exactly as stated:
\[
0<n<10^{10},\qquad n\not\equiv0\pmod{10}.
\]
The repeated-block constructions show that the full sequence is infinite, so the 22 terms form a verified finite initial segment rather than a complete enumeration of all terms.

\section{Conclusion}

The outer-in decimal digit transformation gives a distinct prime-support preservation problem. The reduction
\[
\rad(n)=\rad(T(n))
\iff
\rad\!\left(\frac{n}{g}\frac{T(n)}{g}\right)\mid g,
\qquad
g=\gcd(n,T(n)),
\]
provides an exact characterization and an efficient test for exhaustive computation.

Under the explicit convention that the input does not end in zero, the exhaustive computation finds exactly 22 solutions below \(10^{10}\). The repeated-block argument then gives infinite families beginning with
\[
224\mapsto242,\qquad
448\mapsto484,\qquad
117\mapsto171.
\]

\begin{thebibliography}{9}

\bibitem{OEISA110751}
OEIS Foundation,
\emph{A110751: Numbers n such that n and its digital reversal have the same prime divisors}.
\url{https://oeis.org/A110751}.

\bibitem{OEISA110819}
OEIS Foundation,
\emph{A110819: Non-palindromes in A110751}.
\url{https://oeis.org/A110819}.

\bibitem{OEISA244514}
OEIS Foundation,
\emph{A244514: Numbers n such that the integers formed by all cyclic permutations of the decimal digits of n have the same prime divisors}.
\url{https://oeis.org/A244514}.

\bibitem{Holt2017}
Benjamin V. Holt,
``On Permutiples Having a Fixed Set of Digits,''
\emph{INTEGERS} 17 (2017), A20.
\url{https://arxiv.org/abs/1511.02033}.

\bibitem{Holt2024}
Benjamin V. Holt,
``A Method for Finding All Permutiples with a Fixed Set of Digits from a Single Known Example,''
\emph{INTEGERS} 24 (2024), A88.
\url{https://arxiv.org/abs/2312.02056}.

\end{thebibliography}

\end{document}
